\chapter*{2024年\textsf{\textsc{Hdu}}「高等数学A1」期中试题}

\ExplSyntaxOn
\vspace{-12.0pt}\setcounter{page}{1}
\begin{center}
  \begin{minipage}{.66\textwidth}
    \centering
    \begin{tasks}
      [ label-width = 5\ccwd, item-indent = 2cm,
        after-item-skip = 1ex ](2)
      \task[\textbf{课程信息}：] 理学院/A0714201
      \task[\textbf{解析提供}：] \textsf{\scshape Liu A.N.}
      \task[\textbf{授课教师}：] 大学数学教学团队
      \task[{\makebox[2em][c]{\textbf{\textsf\TeX}}}\textbf{排版}：] \textsf{Izumi~Sakai}
    \end{tasks}
  \end{minipage}
  \begin{minipage}{.15\textwidth}
    \centering
    \begin{tikzpicture}
      \node [ inner~sep = 0 ] (wechat) {
        \fancyqr
          [ size = 4em, image~padding = 1,
            image = \raisebox{4em/33}{\textcolor{DarkGreen \l__whats_darkmode_fading_tl}{\faIcon{weixin}}},
            r~color = DarkGreen \l__whats_darkmode_fading_tl, version = 4
          ]{http://weixin.qq.com/r/sRJiemTEGjllrSAZ90ej}
        };
      \node [ inner~sep = .5ex, align = center, right, font = \scriptsize ]
        at (wechat.east) {未\\央\\学\\社};
    \end{tikzpicture}
  \end{minipage}
  \begin{minipage}{.15\textwidth}
    \centering
    \begin{tikzpicture}
      \node [ inner~sep = 0 ] (qq) {
        \fancyqr
          [ size = 4em, image~x~padding = 2, image~y~padding = 1,
            image = \raisebox{4em/33}{\textcolor{MidnightBlue \l__whats_darkmode_fading_tl}{\faIcon{qq}}},
            l~color = MidnightBlue \l__whats_darkmode_fading_tl, version = 4
          ]{https://qm.qq.com/q/rbfYuqPhQs}
        };
      \node [ inner~sep = .5ex, align = center, left, font = \scriptsize ]
        at (qq.west) {数\\学\\营};
    \end{tikzpicture}
  \end{minipage}
\end{center}
\ExplSyntaxOff

\section{选择题（每题3分，共24分）}

\begin{problem}
  函数 $y = f(x)$ 在 $x = x_0$ 点可微，则在 $x = x_0$ 点
  \begin{tasks}(2)
    \task[\true] 函数 $y = f(x)$ 一定连续
    \task 函数 $y = f(x)$ 可能不连续
    \task 函数 $y = f(x)$ 不一定可导
    \task 函数 $y = f(x)$ 可能是间断点
  \end{tasks}
\end{problem}
\begin{solution}
  可微与可导为等价关系，可微必可导，而可导一定连续. 所以A选项正确.
\end{solution}

\begin{problem}
  极限 $\lim\limits_{x \to \infty} \ab(\frac{x}{3 +  x})^x$ 的结果是
  \begin{tasks}(4)
    \task $e^{-1}$
    \task $e^3$
    \task[\true] $e^{-3}$
    \task $e$
  \end{tasks}
\end{problem}
\begin{solution}
  本题考查重要极限 $\lim\limits_{x \to \infty} \ab(1 + \frac1x)^x = e$.
  \[
    \lim\limits_{x \to \infty} \ab(\frac{x}{3 + x})^x
  = \lim\limits_{x \to \infty} \ab(1 - \frac{3}{x + 3})^x
  = \ab[\lim\limits_{x \to \infty} \ab(1 - \frac{3}{x + 3})^{-(x + 3)/3}]^{-3}
    \cdot \lim\limits_{x \to \infty} \ab(1 - \frac{3}{x + 3})^{-3} = e^{-3}
  \]
\end{solution}

\begin{problem}
  当 $x \to 0$ 时，$f(x) = (1 + x)^\alpha - 1$ 与 $g(x) = \ln(1 + 4x)$
  是等价无穷小，则 $a = $
  \begin{tasks}(4)
    \task $-1$
    \task $1$
    \task $2$
    \task[\true] $4$
  \end{tasks}
\end{problem}
\begin{solution}
  当 $x \to 0$ 时，$\ln (1 + x) \sim x,~ (1 + x)^\alpha - 1 \sim \alpha x$.
  由题意知
  $ \lim\limits_{x \to 0} \frac{(1 + x)^\alpha - 1}{\ln(1 + 4x)}
  = \lim\limits_{x \to 0} \frac{\alpha x}{4x} = 1,~ \alpha = 4$.
\end{solution}

\begin{problem}
  设 $f(x) =
  \begin{cases}
    \frac{x}{\sin x},     & x > 0\\
    0,                    & x = 0\\
    \frac{\arcsin x}{2x}, & x < 0
  \end{cases}$，则 $x = 0$ 是 $f(x)$ 的
  \begin{tasks}(4)
    \task 连续点
    \task[\true] 跳跃间断点
    \task 可去间断点
    \task 无穷间断点
  \end{tasks}
\end{problem}
\begin{solution}
  $ \lim\limits_{x \to 0^+} f(x)
  = \lim\limits_{x \to 0^+} \frac{x}{\sin x} = 1,~
    \lim\limits_{x \to 0^-} f(x)
  = \lim\limits_{x \to 0^-} \frac{\arcsin x}{2x} = \frac12$.
  由于两侧极限 $\lim\limits_{x \to 0^+} f(x)$ 与 $\lim\limits_{x \to 0^-} f(x)$
  均存在，且 $\lim\limits_{x \to 0^+} f(x) \neq \lim\limits_{x \to 0^-} f(x)$，
  所以 $x = 0$ 是 $f(x)$ 的跳跃间断点.
\end{solution}

\begin{problem}
  设在区间 $[0,1]$ 上 $f''(x) > 0$，
  则 $f'(0),~ f'(1) - f(0)$ 或 $f(0) - f(1)$ 的大小关系是
  \begin{tasks}(2)
    \task $f'(1) > f'(0) > f(1) - f(0)$
    \task[\true] $f'(1) > f(1) - f(0) > f'(0)$
    \task $f(1) - f(0) > f'(1) > f'(0)$
    \task $f'(1) > f(0) - f(1) > f'(0)$
  \end{tasks}
\end{problem}
\begin{solution}
  在 $[0,1]$ 上 $f''(x) > 0$，所以 $f'(x)$ 在 $[0,1]$ 上单调递增，$f'(1) > f'(0)$.
  对 $f(x)$ 在 $[0,1]$ 上应用拉格朗日中值定理得 $\exists \xi \in (0,1)$，
  使得 $\frac{f(1) - f(0)}{1 - 0} = f'(\xi)$.
  由于 $f'(1) > f'(\xi) > f'(0)$，所以 $f'(1) > f(1) - f(0) > f'(0)$.
\end{solution}

\begin{problem}
  若函数 $y = f(x)$ 的微分为 $\d y = \sin 2x \cdot e^{\sin^2 x} \d x$，则 $y = $
  \begin{tasks}(4)
    \task $e^{\tan x} + C$
    \task[\true] $e^{\sin^2x} + C$
    \task $\sin x \cdot e^{\cos x} + C$
    \task $e^{\sec x} + C$
  \end{tasks}
\end{problem}
\begin{solution}
  对式子两边同时积分得
  $y = \int \sin 2x \cdot e^{\sin^2x} \d x = e^{\sin^2x} + C$.
\end{solution}

\begin{problem}
  下列选项中是曲线 $y = \frac{1}{x^2 + 3x + 2}$ 的铅直渐进线的有
  \begin{tasks}(4)
    \task[\true] $x = -2$
    \task $y = 1$
    \task $y = 2$
    \task $x = 1$
  \end{tasks}
\end{problem}
\begin{solution}
  $y = \frac{1}{(x + 1)(x + 2)}$，所以 $y$ 有两条铅直渐进线，
  分别是 $x = -1$ 和 $x = -2$.
\end{solution}

\begin{problem}
  使函数 $f(x) = (x^2 - x^4)^{1/3}$ 满足罗尔定理条件的区间为
  \begin{tasks}(4)
    \task $[-2,1]$
    \task $[-1,1]$
    \task[\true] $[0,1]$
    \task $[0,2]$
  \end{tasks}
\end{problem}
\begin{solution}
  $f(-2) \neq f(1),~ f(0) \neq f(2)$. 所以A、D选项错误.
  $ f(x) = (x^2 - x^4)^{1/3},~
    f'(x) = \frac13\frac{2x - 4x^3}{(x^2 - x^4)^{2/3}}$.
  所以在 $x = 0$ 处 $f(x)$ 的导数不存在，B选项错. 所以C选项正确.
\end{solution}

\section{填空题（每题3分，共12分）}

\begin{problem}
  已知 $f(x)$ 可导，函数 $y = f(\arcsin x)$ 的微分
  \result{$(1 - x^2)^{-1/2} \cdot f'(\arcsin x) \d x$}.
\end{problem}
\begin{solution}
  由复合函数求微分法则得
  $\d y = \frac{1}{\sqrt{1 - x^2}} \cdot f'(\arcsin x) \d x$.
\end{solution}

\begin{problem}
  函数 $y = \ln(1 - x)$，则 $y''' = $ \result{$-2(1 - x)^{-3}$}.
\end{problem}
\begin{solution}
  $y' = -(1 - x)^{-1},~ y'' = -(1 - x)^{-2},~ y''' = -2(1 - x)^{-3}$.
\end{solution}

\begin{problem}
  曲线 $y = x^x$ 在点 $(1,1)$ 处的切线方程为 \result{$y = x$}.
\end{problem}
\begin{solution}
  利用对数求导法.
  $\ln y = x \ln x,~ \frac{y'}{y} = \ln x + 1,~ y' = x^x(\ln x + 1)$.
  当 $x = 1$ 时，$y' = 1$. 所以切线方程为 $y - 11 = x - 1$，即 $y = x$.
\end{solution}

\begin{problem}
  已知 $ \lim\limits_{x \to 0} (1 - x)^{a/x}
      = \lim\limits_{x \to 0} \frac{\arctan 2x}{x}$，
  则 $a = $ \result{$-\ln 2$}.
\end{problem}
\begin{solution}
  $ \lim\limits_{x \to 0} \frac{\arctan 2x}{x}
  = 2\lim\limits_{x \to 0} \frac{\arctan 2x}{2x} = 2,~
    \lim\limits_{x \to 0} (1 - x)^{a/x}
  = [\lim\limits_{x \to 0} (1 - x)^{-1/x}]^{-a} = e^{-a} = 2$.
  所以 $a = -\ln 2$.
\end{solution}

\section{计算题A（每题7分，共28分）}

\begin{problem}
  $\lim\limits_{x \to 0} (1 + 3x)^{2/\sin x}$.
\end{problem}
\begin{solution}
  $ \lim\limits_{x \to 0} (1 + 3x)^{2/\sin x}
  = \lim\limits_{x \to 0} \exp\ab[\ln(1 + 3x)^{2/\sin x}]
  = \lim\limits_{x \to 0} \exp\ab[\frac{2}{\sin x} \ln(1 + 3x)]
  = \exp\ab(\lim\limits_{x \to 0} 6x/\sin x) = e^6$.
\end{solution}

\begin{problem}
  求由方程 $xy + e^y + x^2 - 1 = 0$ 所确定的隐函数 $y = y(x)$
  在 $x = 0$ 处的二阶导数 $y''(0)$.
\end{problem}
\begin{solution}\leavevmode
  \begin{itemize}
    \item 方程两边对 $x$ 求导得 $y + xy' + e^y \cdot y' + 2x = 0,~ y'(0) = 0$.
    \item 两边对 $x$ 再求导得
    $y' + y' + xy'' + e^y \cdot y'^2 + e^y \cdot y'' + 2 = 0$.
    将 $y(0) = 0,~ y'(0) = 0$ 代入得到 $y''(0) = -2$.
  \end{itemize}
\end{solution}

\begin{problem}
  $y = \sec x + e^{\sqrt x} + 2024$，求导函数 $\odv y/x$.
\end{problem}
\begin{solution}
  $\odv y/x = \sec x \tan x + \frac12 x^{-1/2}e^{\sqrt x}$.
\end{solution}

\begin{problem}
  求由参数方程 $
  \begin{cases}
    x = t - e^t\\
    y = \cos t
  \end{cases}$
  所确定函数的一阶导 $\odv y/x$ 和二阶导数 $\d^2 y/\d x^2$.
\end{problem}
\begin{solution}
  $ \odv y/x = -\sin t(1 - e^t),~
    \d^2 y/\d x^2 = [e^t(\cos t - \sin t) - \cos t]/(1 - e^t)^3$.
\end{solution}

\section{计算题B（每题8分，共24分）}

\begin{problem}
  设函数 $f(x)$ 在 $x = 1$ 处连续，
  且 $\lim\limits_{x \to 1} \frac{f(x)}{x - 1} = 2$，
  求曲线 $y = f(x)$ 在点 $(1,f(1))$处的切线.
\end{problem}
\scratch
\begin{solution}
  $f(1) = \lim\limits_{x \to 1} f(x)
        = \lim\limits_{x \to 1} \ab[(x - 1)\frac{f(x)}{x - 1}] = 0,~
   f'(1)= \lim\limits_{x \to 1} \frac{f(x) - f(1)}{x - 1} = 2$.
  所以切线为 $y = 2x - 2$.
\end{solution}

\begin{problem}
  证明 $x > 0$ 时，$\frac{x}{1 + x} < \ln(1 + x) < x$.
\end{problem}
\scratch
\begin{solution}
  \begin{proof}
    考察函数 $f(t) = \ln (1 + t)$ 在区间 $[0,x]$ 上满足拉格朗日中值定理
    \[\exists \xi \in (0,x),~ f(x) - f(0) = f'(\xi)(x - 0)\]
    即 $\ln(1 + x) = \frac{x}{1 + \xi}$，其中 $\xi \in (0,x)$. 故
    $\frac{x}{1 + x} < \ln(1 + x) = \frac{x}{1 + \xi} < x$.
  \end{proof}
\end{solution}

\begin{problem}
  写出函数 $f(x) = \ln(1 + x)$ 带有佩亚诺型余项的 $n$ 阶麦克劳林展开式.
\end{problem}
\scratch
\begin{solution}
  将 $f^{(k)}(x) = (-1)^{k - 1} \frac{(k - 1)!}{(1 + x)^k}$ 代入麦克劳林展开式得
  $ \ln (1 + x) = x - \frac{x^2}{2} + \frac{x^3}{3} - \cdots
  + (-1)^{n - 1}\frac{x^n}{n} + o(x^n)$.
\end{solution}

\section{应用题（共12分）}

\begin{problem}[本题8分]
  设 $f(x) =
  \begin{cases}
    x^2 \sin\frac1x, & x \neq 0\\
    a,               & x = 0
  \end{cases}$. 若函数 $f(x)$ 在 $x = 0$连续
  \begin{tasks}[ label = \arabic*. ](2)
    \task 求 $a$ 的值.
    \task 求出函数 $f(x)$ 的导函数 $f'(x)$ 的表达式.
  \end{tasks}
\end{problem}
\scratch
\begin{solution}\leavevmode
  \begin{enumerate}
    \item 函数 $f(x)$ 在 $x = 0$ 连续，故
    $ \lim\limits_{x \to 0} f(x)
    = \lim\limits_{x \to 0} x^2 \sin\frac1x = 0 = f(0) = a$，即 $a = 0$.
    \item 当 $x \neq 0$ 时，$f'(x) = 2x \sin\frac1x - \cos\frac1x$.
    当 $x = 0$ 时，$ f'(0) = \lim\limits_{x \to 0} \frac{f(x) - f(0)}{x - 0}
                          = \lim\limits_{x \to 0} \frac{x^2 \sin1/x}{x} = 0$.
    所以 $f(x)$ 的导函数表达式为
    \[
      f'(x) =
      \begin{cases}
        2x \sin\frac1x - \cos\frac1x, & x \neq 0\\
        0,                            & x = 0
      \end{cases}
    \]
  \end{enumerate}
\end{solution}

\begin{problem}
  已知 $f(x)$ 在区间 $[0,1]$ 连续，在 $(0,1)$ 内可导.
  $f(0) = 0$，证明存在 $\xi \in (0,1)$，使得 $f(\xi) = (1 - \xi)f'(\xi)$.
\end{problem}
\scratch
\begin{solution}
  \begin{proof}
    构造辅助函数 $F(x) = (1 - x)f(x)$ 满足罗尔定理.
    $F(1) = F(0) = 0,~ \exists \xi \in (0,1)$，使得
    $F'(\xi) = -f(\xi) + (1 - \xi)f'(\xi) = 0$. 所以
    $f(\xi) = (1 - \xi)f'(\xi)$.
  \end{proof}
\end{solution}